% ============================================================================
% Chapter 1: Introduction
% ============================================================================

\chapter{Introduction}
\label{chap:introduction}

\section{Overview}
\label{sec:overview}

The field of robotics has witnessed tremendous growth over the past decades, evolving from single-purpose industrial manipulators to sophisticated autonomous systems capable of operating in complex, unstructured environments \cite{siciliano2016robotics}. Among the most promising advancements in this domain is the emergence of \textbf{swarm robotics} --- a paradigm inspired by biological systems such as ant colonies and beehives \cite{dorigo2021swarm}. It emphasizes collective behavior, where simple individual agents interact locally with one another and their environment to achieve complex global objectives.

One of the most critical challenges in autonomous navigation is Simultaneous Localization and Mapping (SLAM) \cite{thrun2005probabilistic}. While single-robot SLAM has matured significantly, it remains limited by its reliance on a single point of failure and its linear progression in time. This project explores the implementation of a collaborative 2D mapping system using robotic swarms. By developing multiple agents equipped with LiDAR and depth cameras, the system can explore indoor environments and detect victims --- such as in office buildings, warehouses, or disaster zones --- concurrently, fusing individual data into a singular, cohesive global map.

\section{Project Motivation}
\label{sec:motivation}

The motivation for developing a swarm-based 2D mapping system is rooted in the transition from centralized, rigid automation to decentralized, flexible intelligence. Beyond basic efficiency, the following factors provide the core impetus for this research:

\begin{figure}[H]
    \centering
    \begin{minipage}[t]{0.48\textwidth}
        \centering
        \includegraphics[width=\linewidth, height=4.5cm]{Project_motivation.png}
        \caption{Project Motivation}
        \label{fig:proj_motivation}
    \end{minipage}
    \hfill
    \begin{minipage}[t]{0.48\textwidth}
        \centering
        \includegraphics[width=\linewidth, height=4.5cm]{Victim_localization.png}
        \caption{Swarm Victim Localization Example}
        \label{fig:victim_loc}
    \end{minipage}
\end{figure}

\subsection{Spatial Parallelism and Temporal Efficiency}
In traditional SLAM, the time required to map an environment is directly proportional to the total area and the complexity of the floor plan. By utilizing a swarm, we implement spatial parallelism. The mapping task is treated as a distributed computing problem where each robot explores a unique ``frontier''. This reduces the total mission time significantly, which is a critical factor in time-sensitive scenarios like gas leak detection or post-earthquake structural assessment.

\subsection{Scalability and Modular Deployment}
One of the most compelling motivations for swarm robotics is scalability. Unlike single-robot systems, where adding capability requires a hardware overhaul, a swarm can be scaled by simply deploying additional low-cost units. This modularity allows the system to adapt to environments of varying sizes --- from a small apartment to a multi-story warehouse --- without needing to rewrite the underlying coordination algorithms. This approach inherently improves cost-effectiveness, as multiple simple robots can achieve tasks that would typically require a single, expensive, and complex robot.

\subsection{Decentralized Intelligence and Reduced Communication Overhead}
Centralized systems suffer from ``bottlenecking,'' where a single central server must process all incoming LiDAR data, leading to latency and potential crashes. Our motivation lies in leveraging decentralized coordination. By processing local maps on-board and only sharing essential ``map fragments'' or ``keyframes'' with neighbors, the swarm minimizes the bandwidth required. This makes the system viable in buildings with thick concrete walls or poor Wi-Fi signals where a central link would fail.

% Add your content here...

\section{Problem Statement}
\label{sec:problem_statement}

Despite the theoretical advantages of swarm robotics, practical implementation faces several challenges:

\begin{enumerate}
    \item \textbf{Lack of Unified Architecture:} Existing systems often use heterogeneous components that lack interoperability.
    \item \textbf{Complex Coordination:} Achieving synchronized behavior among multiple autonomous agents requires sophisticated communication and decision-making protocols.
    \item \textbf{Hardware Constraints:} Designing cost-effective, robust, and modular robot platforms remains challenging.
    \item \textbf{Navigation Complexity:} Multi-robot path planning must account for collision avoidance, traffic management, and dynamic obstacles.
\end{enumerate}

This project addresses these challenges by developing a comprehensive, unified architecture that integrates mechanical design, electrical systems, and autonomous software into a cohesive framework.

% Add your content here...

\section{Project Objectives}
\label{sec:objectives}

The primary objectives of the \textbf{AUSRA} (Autonomous Unified Swarm Robotics Architecture) project are:

\begin{enumerate}
    \item \textbf{Mechanical Design:}
    \begin{itemize}
        \item Design a unified, modular robot chassis suitable for swarm applications.
        \item Perform systematic selection of materials and actuators using decision matrices.
        \item Validate the design through stress analysis and prototyping.
    \end{itemize}
    
    \item \textbf{Electrical System:}
    \begin{itemize}
        \item Develop a robust power distribution and management system.
        \item Integrate sensors (LiDAR, IMU, encoders) for perception and localization.
        \item Design and implement low-level firmware for motor control and sensor interfacing.
    \end{itemize}
    
    \item \textbf{Autonomous Stack:}
    \begin{itemize}
        \item Implement sensor fusion for accurate localization.
        \item Develop path planning algorithms for navigation.
        \item Create a communication framework for swarm coordination.
        \item Design omnidirectional motion control for agile maneuvering.
    \end{itemize}
    
    \item \textbf{Integration and Testing:}
    \begin{itemize}
        \item Integrate all subsystems into a functional prototype.
        \item Conduct comprehensive testing and performance evaluation.
        \item Demonstrate swarm behavior with multiple robots.
    \end{itemize}
\end{enumerate}

% Add your content here...

\section{Thesis Structure}
\label{sec:thesis_structure}

This thesis is organized into seven chapters as follows:

\begin{description}
    \item[Chapter 1: Introduction] Provides an overview of the project motivation, problem statement, objectives, and thesis organization.
    
    \item[Chapter 2: Literature Review] Reviews existing swarm architectures, mechanical design precedents, and control strategies.
    
    \item[Chapter 3: Mechanical Design \& Structural Analysis] Details the design requirements, mechanical trials, actuator selection, hardware assembly, and stress analysis.
    
    \item[Chapter 4: Electrical System \& Low-Level Configuration] Covers the electrical topology, component specifications, PCB design, and firmware implementation.
    
    \item[Chapter 5: High-Level Configuration \& Autonomous Stack] Describes the system architecture, communication stack, localization, planning, and control subsystems.
    
    \item[Chapter 6: Experimental Results \& Discussion] Presents the testing methodology, experimental results, and analysis.
    
    \item[Chapter 7: Conclusion \& Future Recommendations] Summarizes the achievements and suggests directions for future work.
\end{description}
